% Two safe states, what each costs to reach, and how they map to the standards.
\begin{tikzpicture}[font=\sffamily\scriptsize, x=1cm, y=1cm]

\node[chlabel, anchor=west] at (-6.8,3.9) {the same joint, two mountings, two different safe states};

% ================================================================ horizontal
\node[layerlabel] at (-6.8,3.45) {a horizontal or non-backdrivable axis};
\node[linkbody, minimum width=26mm] (h1) at (-5.0,2.15) {};
\node[chlabel, anchor=center] at (-5.0,2.15) {at rest, no energy};
\node[jointpin] (hp) at (-3.5,2.15) {};
\node[linkbody, minimum width=22mm, rotate=8, anchor=west] at (-3.5,2.15) {};
\node[regcell, fill=usrcol, minimum width=54mm, text width=52mm, inner sep=3pt,
      anchor=north west, align=left] at (-6.8,1.5)
  {\textbf{safe state: torque removed.}\\[2pt]
   Immediate, and it \textbf{persists with no energy at all}. Nothing has to
   keep working for it to stay safe, which is the strongest property a safe
   state can have.};

% ================================================================ vertical
\node[layerlabel] at (0.2,3.45) {a vertical axis carrying a payload};
\node[jointpin] (vp) at (2.0,2.45) {};
\node[linkbody, minimum width=22mm, rotate=-70, anchor=west] at (2.0,2.45) {};
\draw[faultedge] (3.0,2.2) -- (3.0,1.4);
\node[chlabel, anchor=west, text=red!60!black] at (3.2,1.8) {it descends};
\node[regcell, fill=red!18, minimum width=54mm, text width=52mm, inner sep=3pt,
      anchor=north west, align=left] at (0.2,1.5)
  {\textbf{torque removed is NOT a safe state here.}\\[2pt]
   The safe state is \textbf{the brake engaged with torque removed}, and it is
   reached through a controlled deceleration first. It persists only as long as
   the brake holds.};

% ================================================================ the mapping
\node[chlabel, anchor=west] at (-6.8,-0.85) {and how both map onto the published scales, which is the cleanest fact in this area};

\node[regcell, fill=white, minimum width=114mm, anchor=north west,
      align=left, inner sep=3pt] at (-6.8,-1.25)
  {\begin{tabular}{@{}p{22mm}p{32mm}p{26mm}p{26mm}@{}}
   \textbf{stop category} & \textbf{what it does to actuator power} &
   \textbf{drive sub-function} & \textbf{used here for} \\[1pt]
   \hline
   0 & immediate removal, an uncontrolled stop & safe torque off &
   the horizontal case, demonstrated \\[2pt]
   1 & a controlled stop with power available, then removal once stopped &
   safe stop 1 & the vertical case, in logic only \\[2pt]
   2 & a controlled stop with power remaining available & safe stop 2 &
   not used: this node has no reason to hold position under power \\
   \end{tabular}};

\node[umlnote, text width=54mm, anchor=north west] at (-6.8,-5.0)
  {The category definitions themselves sit behind a paywall. The wording is
   identical across the sources this volume could read, so it is carried here as
   \textbf{attested rather than primary verified}, and the chapter marks it that
   way rather than quietly promoting it.};
\node[umlnote, text width=54mm, anchor=north west] at (0.2,-5.0)
  {\textbf{There is no stop category 3.} The scale has three values and the
   third one is 2. The frequent appearance of a fourth in requirement documents
   comes from confusing this scale with a reliability scale that runs to four,
   which is a different axis entirely.};
\end{tikzpicture}
